\begin{abstract}
Deep reinforcement learning (DRL) is increasingly used for microgrid battery dispatch, yet policies are often validated mainly by cumulative economic cost. This scalar metric can obscure state-of-charge trajectory shape, raw-action feasibility, and reliance on safety-layer intervention. This paper presents a morphology-aware validation framework that combines Shield-Aware Intervention Learning with Soft Actor-Critic (SAIL-SAC) and a multi-gate checkpoint-selection protocol, P4 BalancedGate. The protocol evaluates policies using mid-band state-of-charge dwelling, price-responsive cycling, infeasible-action suppression, and logged intervention metrics. In an IEEE 33-bus distribution-network stress case with 2023 training data and full-year 2024 evaluation, the six P4-selected candidates record mean objective cost $11.788 \pm 0.120$\,M\textyen{} (million CNY-equivalent), $20.4 \pm 15.6$\% mid-band dwelling, and $6.1 \pm 7.3$\% infeasible-action dwell under the simple efficiency-based model. Cross-fidelity replay shows small objective variation across battery representations, with mean cost changing from 11.788\,M\textyen{} under the simple EBM to 11.737\,M\textyen{} under the Thevenin full PBM. The full-year evaluation records zero shield activation but a nonzero normalized inventory-teacher action gap ($0.301 \pm 0.087$), showing that intervention diagnostics capture behavior not visible in cost alone. The results indicate that morphology and intervention diagnostics can serve as complementary criteria when selecting interpretable, constraint-aware battery-dispatch policies.
\end{abstract}

\begin{keyword}
Battery energy storage system\sep distribution-network microgrid\sep shielded reinforcement learning\sep morphology-aware validation\sep soft actor--critic\sep intervention diagnostics
\end{keyword}
